\section{Exterior powers and Adams operations with every split-support label retained}
\label{sec:full-lattice-adams}

The human source for the Adams-to-arithmetic-topos construction is
Alain Connes and Caterina Consani,
\href{https://arxiv.org/abs/2004.08879v1}{\emph{Segal's Gamma rings and
universal arithmetic}}, original author source \nolinkurl{Froblifts.tex}.
The source ranges read for the present calculation are 600--704,
with the Adams formulas and their full topos argument at 632--681,
and the Atiyah--Tall bibliography entry at 1242--1243. The source
SHA--256 is
\par\noindent
\nolinkurl{226eecb4cfde540362b1d5073fb23bad2c3de6e0b8c3cf267b47447e3166298d}.
\par\noindent
Their mathematical reference is M.~F.~Atiyah and D.~O.~Tall,
\emph{Group representations, $\lambda$-rings and the
$J$-homomorphism}, Topology \textbf{8} (1969), 253--297,
especially their Section 5 and Proposition 5.2 as identified in
Connes--Consani's text. No claim to have read that entire reference
is made here. The twistor maps and real structure are Connes--Consani's
\href{https://arxiv.org/abs/2609.00299v1}{\emph{The Absolute Twistor
Line and the Geometry of $\overline{\operatorname{Spec}\mathbb Z}$}};
the original \nolinkurl{CC.tex} was read in full for the preceding
HA1--29 derivation. Its ordinary $K_0$ calculation, precise equivariant
correction, and real-lift maps are HA23--29 of the retained proof
\nolinkurl{SIGNED_HOPF_AND_ADAMS.tex}. The operations below are a
full-support extension derived here, not an assertion attributed
to either human source.

\subsection{The original ring and its integral exterior-power formula}

Retain the original generator, without changing its sign:
\begin{equation}\tag{FLA1}
 B=K_0(\mathbb P^1_{\mathbb C})
       =\mathbb Z[\eta]/(\eta^2),\qquad
 \eta=[\mathcal O(1)]-1,
 \qquad b=r+a\eta\quad(r,a\in\mathbb Z).
\end{equation}
The presentation and uniqueness of these coordinates are proved in
HA23--25 by the Euler sequence and the diagonal resolution. In
particular $(r,a)$ here is rank and degree, including negative
virtual rank and degree. Let $\ell=[\mathcal O(1)]=1+\eta$.
The ordinary exterior-power operations give
$\lambda_t(1)=1+t$ and $\lambda_t(\ell)=1+\ell t$.
For an exact sequence of vector bundles, the exterior powers have
the filtration with associated graded terms
$\bigwedge^i E'\otimes\bigwedge^{n-i}E''$; it is obtained on
locally split charts and is independent of the splitting. Taking
classes proves multiplicativity of total exterior power under
addition in $K_0$. Extension to virtual differences uses the formal
inverse of a series with constant coefficient one. Since
$b=(r-a)1+a\ell$, this proves the exact integral formula
\begin{equation}\tag{FLA2}
 \lambda_t^B(r+a\eta)
  =(1+t)^{r-a}(1+(1+\eta)t)^a
  =(1+t)^r\left(1+\frac{a\eta t}{1+t}\right).
\end{equation}
Every inverse here is a formal power-series inverse over $B$,
not division by an integer. In particular $1/(1+t)=\sum_{q\ge0}
(-t)^q$. For every integer $a$, the identity $(1+v)^a=1+av$
holds when $v^2=0$: positive exponents follow by multiplication,
and negative ones follow from $(1+v)^{-1}=1-v$. This proves the
second equality for all the stated coordinates.

For $r\in\mathbb Z$ and $n\ge0$, define the integer
$\binom rn$ by the coefficients of $(1+t)^r$. Equivalently it
is the usual falling-product formula, with integrality for negative
$r$ given by $\binom rn=(-1)^n\binom{n-r-1}{n}$. Thus (FLA2)
gives
\begin{equation}\tag{FLA3}
 \lambda_0^B(b)=1,\qquad
 \lambda_n^B(r+a\eta)=
 \binom rn+a\binom{r-1}{n-1}\eta\quad(n\ge1).
\end{equation}
The expressions remain valid for $r=0$, $a=0$, or either coordinate
negative. Direct multiplication in (FLA2), with $\eta^2=0$, proves
\begin{equation}\tag{FLA4}
 \lambda_t^B(b+c)=\lambda_t^B(b)\lambda_t^B(c)
                    \qquad(b,c\in B).
\end{equation}
This verifies the required identity without omitting the virtual
classes.

\subsection{The support-preserving operations and every zero fibre}

Let $L$ be any nontrivial bounded distributive lattice, finite or
infinite. Keep the original semiring
\begin{equation}\tag{FLA5}
 \begin{gathered}
 S=G_L(B)=\{(0,\lambda):\lambda\in L\}
                  \cup\{(b,1_L):b\in B\},\\
 (b,\lambda)+(c,\mu)=(b+c,\lambda\vee\mu),\qquad
 (b,\lambda)(c,\mu)=(bc,\lambda\wedge\mu),\\
 z_\lambda=(0,\lambda),\quad \tau=z_{0_L},\quad e=z_{1_L},
 \quad 1_S=(1,1_L),\quad \overline{(b,\lambda)}=(-b,\lambda).
 \end{gathered}
\end{equation}
The bar is the actual support-preserving signed operation; it is
not a global additive inverse. In particular
$x+\overline x=z_\lambda$ when $x=(b,\lambda)$.
Define
\begin{equation}\tag{FLA6}
 \lambda_0^{\uparrow}(x)=1_S,\qquad
 \lambda_n^{\uparrow}(b,\lambda)
        =(\lambda_n^B(b),\lambda)\quad(n\ge1),\qquad
 \Lambda_t(x)=\sum_{n\ge0}\lambda_n^{\uparrow}(x)t^n.
\end{equation}
These are defined on the original carrier: if $\lambda\ne1_L$
then $b=0$ and $\lambda_n^B(0)=0$ for $n\ge1$. The coefficient
at every positive degree therefore has exactly the stated support,
even when its amplitude vanishes. In particular
\begin{equation}\tag{FLA7}
 \begin{gathered}
 \lambda_1^{\uparrow}(x)=x,\qquad
 \lambda_n^{\uparrow}(\tau)=\tau,\qquad
 \lambda_n^{\uparrow}(z_\lambda)=z_\lambda\quad(n\ge1),\\
 \lambda_1^{\uparrow}(1_S)=1_S,\qquad
 \lambda_n^{\uparrow}(1_S)=e\quad(n\ge2).
 \end{gathered}
\end{equation}
In the formal semiring $S[[t]]$, the constant series $1$ has zero
coefficient $\tau$ at positive degrees. Accordingly (FLA7) does
not identify $\Lambda_t(1_S)$ with the polynomial $1+t$.

The full additive exterior-power law is nevertheless exact:
\begin{equation}\tag{FLA8}
 \Lambda_t(x+y)=\Lambda_t(x)\Lambda_t(y),\qquad
 \Lambda_t(\tau)=1.
\end{equation}
To check all labels, write $x=(b,\lambda)$ and $y=(c,\mu)$.
For degree $n\ge1$, the terms at indices $(0,n)$ and $(n,0)$
on the right have supports $\mu$ and $\lambda$. Every other
term has support $\lambda\wedge\mu$. Their join is exactly
$\lambda\vee\mu$, the support on the left. Their amplitudes
agree by (FLA4). At degree zero both sides are $1_S$. This proves
the identity in $S[[t]]$, not only after applying an amplitude map.
Thus (FLA6)--(FLA8) specify the support-preserving exterior-power
laws used here. They do not assert that $S$ is a ring or that the
ordinary special-$\lambda$-ring axioms hold on it: already the
ordinary unit axiom $\lambda_n(1)=0$ for $n>1$ would replace
$e$ in (FLA7) by the different element $\tau$.

\subsection{The exact relative inverses and their support domains}

Define the idempotent series
\begin{equation}\tag{FLA9}
 E_\lambda(t)=1+\sum_{n\ge1}z_\lambda t^n
       =\Lambda_t(z_\lambda),\qquad
 E_\lambda E_\mu=E_{\lambda\vee\mu},\qquad E_{0_L}=1.
\end{equation}
The middle identity follows by the same support computation as
(FLA8), with all positive-degree amplitudes zero. In particular
$E_\lambda^2=E_\lambda$, and $E_\lambda=1$ holds exactly when
$\lambda=0_L$, as seen in its coefficient at degree one.

The precise multiplicative domains are
\begin{equation}\tag{FLA10}
 \mathcal H_\lambda=
 \left\{1+\sum_{n\ge1}(b_n,\lambda)t^n:
                   (b_n,\lambda)\in S\text{ for every }n\right\}.
\end{equation}
For $\lambda\ne1_L$ this is the one-element group $\{E_\lambda\}$.
For $\lambda=1_L$, the amplitude map identifies $\mathcal H_{1_L}$,
with its own identity $E_{1_L}$, with the ordinary group
$1+tB[[t]]$. To verify the latter without any assumed inverse,
for $A(t)=1+\sum b_nt^n$ set $c_0=1$ and recursively
$c_n=-\sum_{i=1}^n b_i c_{n-i}$. The coefficients are in $B$
and multiplication gives $A(t)\sum c_nt^n=1$. Lifting every
positive coefficient to support $1_L$ makes the product
$E_{1_L}$. The amplitude coordinates also prove uniqueness of
this inverse. This proves the group assertions in every case.
Multiplication of two arbitrary domains satisfies
\begin{equation}\tag{FLA11}
 \mathcal H_\lambda\mathcal H_\mu
           \subseteq\mathcal H_{\lambda\vee\mu}.
\end{equation}
When $\lambda\le\nu$, the map $A\mapsto E_\nu A$ is a group
homomorphism $\mathcal H_\lambda\to\mathcal H_\nu$, sends the
identity to $E_\nu$, and simply raises every positive-degree
support to $\nu$. These maps compose for nested supports, by
(FLA9). Thus all multiplication domains and their connecting maps
are specified.

For $x=(b,\lambda)$ the actual signed inverse obeys
\begin{equation}\tag{FLA12}
 \Lambda_t(x)\Lambda_t(\overline x)=E_\lambda,
 \qquad E_\lambda\Lambda_t(x)=\Lambda_t(x).
\end{equation}
The first follows from (FLA8) and $x+\overline x=z_\lambda$;
the second follows directly by multiplying amplitudes and labels.
Consequently $\Lambda_t(\overline x)$ is exactly the group
inverse in $\mathcal H_\lambda$, and both reflexive inverse laws
$A\overline A A=A$, $\overline A A\overline A=\overline A$
hold in the full series semiring. Except at $\tau$, it cannot be
an ordinary inverse with product $1$: the coefficient of degree
one in a product of $\Lambda_t(x)$ with any series of constant
coefficient $1_S$ has support at least $\lambda$. A global inverse
would have that constant coefficient, since the only possible
constant inverse of $1_S$ is $1_S$ itself. For $\lambda\ne0_L$
this precludes a zero coefficient $\tau$ at degree one.
Finally $\Lambda_t$ is injective because its first coefficient
is $x$. Equations (FLA8)--(FLA12) therefore describe the whole
additive split-support object by an exact series embedding, with
the correct relative, rather than global, inverses.

\subsection{Adams endomorphisms and the division-free Newton calculation}

For every integer $n\ge1$ define
\begin{equation}\tag{FLA13}
 \psi_n^B(r+a\eta)=r+na\eta,\qquad
 \Psi_n(r+a\eta,\lambda)=(r+na\eta,\lambda).
\end{equation}
The lifted map is defined on all of $S$ for the same carrier reason
as (FLA6). It is a unital semiring endomorphism. Additivity follows
from the displayed formula and the unchanged join of the two
supports. For multiplication, in the original coordinates,
\[
 (r+a\eta)(s+b\eta)=rs+(rb+sa)\eta,
\]
and both $\psi_n^B$ of this product and the product of the two
images equal $rs+n(rb+sa)\eta$. The support is the unchanged
meet. Zero and unit are preserved. The same calculation proves
\begin{equation}\tag{FLA14}
 \Psi_1=\operatorname{id},\quad
 \Psi_m\Psi_n=\Psi_{mn},\quad
 \Psi_m\lambda_n^{\uparrow}
       =\lambda_n^{\uparrow}\Psi_m\quad(n\ge0),\quad
 \Psi_m(z_\lambda)=z_\lambda.
\end{equation}
For the commutation with $\lambda_n$, use (FLA3): both maps keep
the binomial rank term and replace its $\eta$ coefficient by $m$
times that coefficient, with the same support. The case $n=0$
is the preserved unit. Thus these operations define an actual
$\mathbb N^\times$-action by semiring endomorphisms, and hence
the corresponding semiring object in the presheaf topos underlying
the arithmetic site. This is the semiring version of the exact
action mechanism in Connes--Consani's Lemma \texttt{lambda1},
source lines 649--655, proved here on the full original carrier.

No division by a positive integer is needed to identify these as
the Newton operations of (FLA6). Formal differentiation of (FLA2)
in the ring $B[[t]]$ gives the integral identity
\begin{equation}\tag{FLA15}
 t\frac{d}{dt}\lambda_t^B(b)\,
                   \lambda_t^B(-b)
  =\frac{rt}{1+t}+\frac{a\eta t}{(1+t)^2}
  =\sum_{n\ge1}(-1)^{n-1}(r+na\eta)t^n.
\end{equation}
For the first equality, differentiate the integer powers in
(FLA2), including their formal inverses, and use $\eta^2=0$.
For the second, differentiate the integral geometric series for
$(1+t)^{-1}$ or multiply it by itself. Its squared inverse has
coefficient $(-1)^q(q+1)$ at degree $q$. All the operations take
place over $\mathbb Z$ and all denominators displayed are units
of $\mathbb Z[[t]]$, so no rationalization is involved.

For an element of support $\lambda$, define a signed multiple in
its additive support fibre by
\[
 [c]_\lambda(b,\lambda)=(cb,\lambda)\qquad(c\in\mathbb Z).
\]
For $c>0$ this is repeated addition; $[-1]_\lambda x=\overline x$.
For $c=0$ it is the identity $z_\lambda$ of that fibre, not an
instruction to change it to the global zero $\tau$. This convention
records exactly the support of signed polynomial terms.
Let
\begin{equation}\tag{FLA16}
 W_x(t)=\sum_{n\ge1}[-1]^{n-1}_\lambda\Psi_n(x)t^n
             \qquad(x=(b,\lambda)),
\end{equation}
where the superscript denotes the iterated signed operation and
the constant coefficient is $\tau$. Every positive coefficient
has support $\lambda$. If $D$ denotes the formal derivative in
the semiring, with repeated-addition coefficient $n$, then the
exact lifted identities are
\begin{equation}\tag{FLA17}
 tD\Lambda_t(x)=\Lambda_t(x)W_x(t),\qquad
 W_x(t)=\Lambda_t(\overline x)\,tD\Lambda_t(x).
\end{equation}
For the first identity, every positive-degree coefficient on each
side has support $\lambda$, and the amplitudes agree by (FLA15)
after multiplying there by $\lambda_t^B(b)$. Both constant
coefficients are $\tau$. For the second, multiply by the exact
relative inverse in (FLA12). Its extra factor $E_\lambda$ acts
identically on $W_x$: its constant coefficient is $1_S$, its
positive amplitudes vanish, and all resulting supports remain
$\lambda$. This proves the second identity without treating
$E_\lambda$ as the constant series $1$.

Taking degree $n\ge1$ in the first identity gives the complete
division-free Newton formula in $S$:
\begin{equation}\tag{FLA18}
 n\lambda_n^{\uparrow}(x)
   =\sum_{i=1}^n[-1]^{i-1}_\lambda
       \bigl(\lambda_{n-i}^{\uparrow}(x)\Psi_i(x)\bigr).
\end{equation}
Every term has support $\lambda$, including $i=n$, since
$\lambda_0^{\uparrow}(x)=1_S$ has top support. The sum is
therefore in that same additive fibre. Each such fibre is a group
under addition with identity $z_\lambda$: it is $B$ at top
support and a singleton at every other support. Thus solving for
the last term uses only a coefficient $+1$ or $-1$ and gives
\begin{equation}\tag{FLA19}
 \Psi_n(x)=[-1]^{n-1}_\lambda
                  \bigl(n\lambda_n^{\uparrow}(x)\bigr)
  +\sum_{i=1}^{n-1}[-1]^{n+i-1}_\lambda
       \bigl(\lambda_{n-i}^{\uparrow}(x)\Psi_i(x)\bigr).
\end{equation}
For $n=1$ the sum is absent and the identity is $\Psi_1=x$.
This recursive formula defines a unique sequence in the stated
support fibres. It proves that (FLA13) is the actual Newton
sequence of the lifted exterior operations. It does not attempt
to reconstruct exterior operations by dividing $n$ in (FLA18).

There is a local sign correction to the human source's displayed
equation \texttt{lognewton}, line 645. It prints $(-1)^n$ on the
right, while lines 637--640 state $\psi_1=\lambda_1$ and the usual
Newton polynomials. With the source's displayed derivative
$t\partial_t\lambda(t)/\lambda(t)$, its coefficient of $t$ is
$\lambda_1$ and therefore requires $(-1)^{n-1}$. For the ordinary
unit class $b=1$, (FLA15) gives $t/(1+t)=t-t^2+t^3-\cdots$;
the printed sign would give its negative. Equations (FLA15)--(FLA19)
prove the corrected formula directly. This correction is to that
display and does not negate the source's topos-action argument.

\subsection{Actual characteristic-prime Frobenius lifts with support retained}

Let $p$ be a rational prime and use the exact quotient map
\begin{equation}\tag{FLA20}
 q_p:S\longrightarrow G_L(B/pB),\qquad
 (b,\lambda)\longmapsto(b\bmod pB,\lambda).
\end{equation}
It preserves both operations and is surjective: every nonzero
amplitude class has a representative at top support, while every
zero label already has its designated representative. Its precise
kernel congruence is equality of supports together with amplitude
difference in $pB$. Thus it does not identify distinct zero labels.
In $B$, for $b=r+a\eta$,
\begin{equation}\tag{FLA21}
 b^p=r^p+p r^{p-1}a\eta,\qquad
 \psi_p^B(b)-b^p=(r-r^p)+p a(1-r^{p-1})\eta\in pB.
\end{equation}
The inclusion uses $r^p\equiv r\pmod p$ for every integer $r$.
For a direct proof, the intermediate coefficients of $(x+1)^p$
are divisible by $p$, since $p$ divides its numerator factorial
and neither denominator factorial. Induction from $0$ gives the
congruence for nonnegative integers; applying the same polynomial
identity backward gives it for negative integers as well.
Since the support of the $p$th power is $\lambda\wedge\cdots
\wedge\lambda=\lambda$, (FLA21) proves the full equality
\begin{equation}\tag{FLA22}
 q_p\Psi_p=\operatorname{Fr}_p q_p,\qquad
 \operatorname{Fr}_p(c,\lambda)=(c^p,\lambda)
                        \text{ on }G_L(B/pB).
\end{equation}
The target map is a semiring endomorphism: its amplitude is the
characteristic-$p$ Frobenius, and its labels preserve joins and
meets. Thus (FLA22) is an actual Frobenius-lift statement on the
full split-support base, not only a congruence in its ring shadow.

\subsection{Odd twistor pullbacks, with their real and equivariant data}

For positive odd $n$, put $\chi(n)=1$ for $n\equiv1\pmod4$
and $\chi(n)=-1$ for $n\equiv3\pmod4$. The original twistor maps
are exactly
\begin{equation}\tag{FLA23}
 F_n(z)=
 \begin{cases}z^n,&n\equiv1\pmod4,\\-1/z^n,&n\equiv3\pmod4.
 \end{cases}
\end{equation}
In homogeneous coordinates they are $[X^n:Y^n]$ and
$[-Y^n:X^n]$, respectively. Both pairs are sections of
$\mathcal O(n)$ with no common zero, proving
$F_n^*\mathcal O(1)\simeq\mathcal O(n)$ and hence
$F_n^*\eta=n\eta$. These two formulas therefore prove on every
original class, and then on its entire support fibre, that
\begin{equation}\tag{FLA24}
 G_L(F_n^*)=\Psi_n,
 \qquad F_n^*(r+a\eta)=r+na\eta.
\end{equation}
The maps compose as $F_mF_n=F_{mn}$ for odd indices: the involution
$a(z)=-1/z$ commutes with odd powering, $a^2=\operatorname{id}$,
and the parity selecting $a$ adds because $\chi(mn)=\chi(m)\chi(n)$.
This proves the geometric counterpart of (FLA14) for these maps.

For the scaling group, however, the exact equivariance is
$F_n(u z)=u^{n\chi(n)}F_n(z)$. Thus a bundle with pole weight
multisets $(W_0,W_\infty)$ pulls back with the specific weights
\begin{equation}\tag{FLA25}
 (W'_0,W'_\infty)=
 \begin{cases}
 (nW_0,nW_\infty),&n\equiv1\pmod4,\\
 (-nW_\infty,-nW_0),&n\equiv3\pmod4.
 \end{cases}
\end{equation}
Indeed the action on $(z,v)$ in the pulled-back fibre is
$u\cdot(z,v)=(uz,u^{n\chi(n)}\cdot v)$; its target fibre is
correct by the displayed equivariance. The second case exchanges
the two poles, proving all labels in (FLA25). The trivial line
with fibre character $u$ then has pullback character $u^{-n}$
in the second case, while its equivariant Adams operation has
character $u^n$. Restriction to the fixed point detects distinct
Laurent monomials, so (FLA24) does not assert their equality
equivariantly. The exact connecting map is the forgetful map
$\operatorname{For}:K_0^{\mathbb C^\times}(\mathbb P^1)
\to B$: forgetting the specified pulled-back action gives ordinary
pullback, hence
\begin{equation}\tag{FLA26}
 G_L(\operatorname{For})\,G_L(P_n)
       =\Psi_n\,G_L(\operatorname{For}),
\end{equation}
where $P_n$ is the equivariant pullback just constructed.
The equality follows on actual bundles and is preserved under
exact-sequence relations, tensor product, and every support label.

For the antipodal real map $\sigma(z)=-1/\overline z$ one has
$F_n\sigma=\sigma F_n$, by the same odd-power calculation and
real coefficients in (FLA23). If an actual bundle has an
antilinear lift $\tau:E_z\to E_{\sigma(z)}$, its transported
lift is
\begin{equation}\tag{FLA27}
 \tau_n(z,v)=(\sigma(z),\tau v),\qquad
 \tau_n^2(z,v)=(z,\tau^2v).
\end{equation}
Commutation of the base maps proves its target fibre, and applying
it twice proves the square. It therefore preserves a supplied
real or quaternionic sign exactly. It does not construct a lift,
filtration, polarization, or connection from an ordinary $K_0$
class. Equations (FLA25)--(FLA27) give the actual maps on the
additional structures that are present.

\subsection{Naturality in the entire support lattice}

Let $\phi:L\to L'$ preserve zero, unit, join, and meet, with no
finiteness requirement. The original morphism is
$G_\phi(b,\lambda)=(b,\phi(\lambda))$. Its well-definedness and
operation laws follow directly from (FLA5); in particular nonzero
amplitudes remain at top support. The explicit formulas prove
\begin{equation}\tag{FLA28}
 \begin{gathered}
 G_\phi\lambda_n^{\uparrow,L}
       =\lambda_n^{\uparrow,L'}G_\phi\quad(n\ge0),\qquad
 G_\phi\Psi_n^L=\Psi_n^{L'}G_\phi,\\
 G_\phi(\overline x)=\overline{G_\phi(x)},\qquad
 G_\phi(E_\lambda)=E_{\phi(\lambda)},\qquad
 G_\phi(\mathcal H_\lambda)\subseteq
                         \mathcal H_{\phi(\lambda)}.
 \end{gathered}
\end{equation}
The series maps are coefficientwise; their constant unit is
preserved. Thus they preserve the relative inverses with their
actual changed identity and carry both sides of (FLA17)--(FLA19)
to the corresponding full identities over $L'$. Reduction modulo
$p$, the ordinary twistor pullback, and its forgetful square commute
with these same maps because they change only the amplitude.
Composition of lattice maps is exact on each coordinate, proving
composition of all these squares. Finally the amplitude and support
maps themselves satisfy
\begin{equation}\tag{FLA29}
 q\Psi_n=\psi_n^Bq,\quad q\lambda_n^{\uparrow}=\lambda_n^Bq,
 \qquad \sigma\Psi_n=\sigma,\quad
 \sigma\lambda_n^{\uparrow}=\sigma\quad(n\ge1).
\end{equation}
Here $q(b,\lambda)=b$ and $\sigma(b,\lambda)=\lambda$; the
constant operation has separately specified top support. This
records the exact relationship of both shadows to the retained
full operations. No cohomological or purity identification is
assumed or inferred.
